\section{Is the constant an artifact of the proof?}
Take $X=\RR$, $T(x)=qx$, $y_0=0$, and add exactly $\eta$ at each update. The recurrence has the explicit solution
\[
 y_n=\eta\sum_{k=0}^{n-1}q^k
     =\eta\frac{1-q^n}{1-q}.
\]
This follows either from induction or by substituting the formula into $y_{n+1}=qy_n+\eta$.

For the induction step, assume $y_n=\eta(1-q^n)/(1-q)$. Then
\begin{align*}
 y_{n+1}
 &=q\,\frac{\eta(1-q^n)}{1-q}+\eta\\
 &=\frac{\eta(q-q^{n+1})+\eta(1-q)}{1-q}\\
 &=\frac{\eta(1-q^{n+1})}{1-q}.
\end{align*}
At $n=0$ the formula gives zero, as required. The exact formula is therefore established for every finite step; the limiting value follows afterward from $q^n\to0$. We do not infer the formula merely from the appearance of the first few terms. Since $x_*=0$ and $y_n\ge0$, the uniform-error bound is an equality. Thus its limiting factor $1/(1-q)$ cannot be decreased for all systems satisfying these assumptions.

A sharp bound need not predict convergence of every allowed approximate sequence. For example, let $T(x)=0$ and choose $y_n=(-1)^n\eta$ with $\eta>0$. This map satisfies the contraction inequality for any chosen $q\in(0,1)$, and
\[
 d(y_{n+1},T(y_n))=\eta.
\]
Yet $y_n$ alternates between two different values and does not converge. Bounded errors give a distance estimate, not convergence by themselves. Vanishing errors provide the extra conclusion in Theorem~\ref{thm:inexact}.
